\section{Scope, theorem inputs, and the central reductions}

The main unrestricted tasks remain open in this manuscript: an explicit polynomial-size hitting set for general width-$w$ ROABPs, a uniformly efficient positive rule for all rectangular and plethysm multiplicities, and a determinant-versus-padded-permanent separation in the required asymptotic regime. Counterexamples below refute specified transfers, not those open problems. Every stated deduction has an identifier. The accompanying event register distinguishes complete proof readings, finite exact computations, dependency checks, and carried background. No Lean reconstruction or formal verification of the new mathematics is claimed.

The principal positive results are exact rather than heuristic. R-01--R-09 identify the CUE character and ratios operations, positive projector formulas, orbit-closure Gram ranks, and the missing quotient numerator. R-10 closes the height-two rectangle and quadratic determinant case. G-05 excludes every at-most-three-row multiplicity obstruction. G-07 uses the reduced determinantal-quartic divisor to give an exact shifted-plethysm formula. R-12 combines it with a certified permanent projection to give an exact four-row permanent-minus-determinant difference. R-13 gives a finite, completely specified Jacobian classification of quaternary permanent dominance. G-08--G-09 turn uniform determinantal constructions into explicit bounds on the number of rows an obstruction must use at a prescribed padding size.

On isolation, I-01 gives the genuine Hoza transfer for independent vertex variables. I-03--I-04 identify the coefficient flags and projected matroid cones required by basis isolation. I-14 gives the weaker joint-compression criterion; this is potentially sufficient even when no individual weight is basis isolating. The subsequent decoder, finite-field, and specialization results test specific repairs. The purity section gives explicit sparse constructions, a sheaf-theoretic Gram criterion with conductor accounting, and precise limitations of energy approximation. P-11 shows that an explicit common number-field evaluation point of polynomial degree already yields an unconditional hitting set by polynomial interpolation; merely assigning that encoding to Chebotarev does not supply the missing construction.

\subsection{Stipulated analytic inputs and supplied material}

We accept Alp\"oge's RH and Dirichlet GRH as theorem inputs, exactly as requested. In \texttt{ALPOGE\_Main.lean}, the exported declarations are \texttt{riemannHypothesis}, \texttt{riemannZeta\_eq\_zero\_re}, \texttt{LFunction\_ne\_zero\_of\_half\_lt\_re}, \texttt{LFunction\_eq\_zero\_re}, and \texttt{LFunction\_eq\_zero\_re\_of\_isPrimitive}; the predicates are defined in \texttt{ALPOGE\_Statements.lean}. The attached note supplies narrative descriptions of its numbered steps. Neither the imported Lean dependency tree nor the claimed all-files scan is in this attachment, and neither was rebuilt in this run. This is an input convention, not a review of the theorem's truth.

Masullo's general number-field Hecke GRH is used only conditionally, as the prompt's attribution to section~1.2. The supplied statement index has no numbered section~1.2 theorem or its proof. Accordingly we give no invented locator. Its Theorem~13.3.1.1, indexed PDF page~522, is a trace-formula claim; Conjecture~14.1.2.1, indexed page~529, concerns automorphic $\mathrm{GL}_n$ GRH and simplicity and remains explicitly conjectural there. These are different statements. The zero-statistics section names the precise hypotheses discharged by the stipulated input and those that remain.

The archive was unpacked once. The source manifest records the actual supplied files and hashes. The five cote~133 files are LaTeX transcriptions, not original page images. Assertions in the accompanying earlier study about visual readings remain credited to that study; our transcription table records mathematical necessity separately from paleographic confidence.

\subsection{Why the three multiplicity problems are distinct}

Let $V=\mathbb C^{n^2}$, $W=\operatorname{Sym}^nV^*$, and $A=\operatorname{Sym}(W^*)$. For a homogeneous form $f\in W$, write
\[
 m_f(\lambda,d)=\operatorname{mult}_{S_\lambda V}\mathbb C[\overline{GL(V)f}]_d,
 \qquad a_\lambda(d[n])=\operatorname{mult}_{S_\lambda V}A_d.
\]
The symbol $g(\lambda,\mu,\nu)$ denotes the ordinary Kronecker coefficient. One must distinguish the ambient plethysm multiplicity, the connected-stabilizer rectangular Kronecker multiplicity, and the coordinate-ring multiplicity of an orbit closure. The transpose component of the determinant stabilizer and the orbit boundary are separate restrictions. R-03 gives finite counterexamples to identifying them. R-05 computes the closure multiplicity by an actual rank, and R-08 expresses the same requirement as a quotient numerator.

If $X\subseteq W$ is $GL(V)$-stable and $V_k$ is spanned by the first $k$ basis vectors, let $X_k$ be the closure of its restriction image in $\operatorname{Sym}^nV_k^*$. For $\ell(\lambda)\leq k$, the multiplicities in $\mathbb C[X]$ and $\mathbb C[X_k]$ agree. The short verification used repeatedly below is as follows. A highest-weight polynomial of weight $(\lambda_1,\ldots,\lambda_k,0,\ldots,0)$ uses only coefficient variables supported on $V_k$, since all generator weights are nonnegative. It vanishes on $X$ exactly when it vanishes on $X_k$. Taking the quotient of the common ambient highest-weight space proves the assertion. This is a coefficient-subalgebra version of the usual inheritance principle, proved here in precisely the form needed for the orbit restrictions.

The exact identities alone impose no efficient complexity bound. In particular, a nonnegative squared-norm formula may have exponentially many terms; a rank formula may require exponentially large matrices; and a positive dimension may become zero after highest-weight projection. The size and sign statements in each result are therefore part of the theorem, not suppressed implementation details.

\subsection{How the specific new analytic tools enter}

Alp\"oge's step~2 (p.~3) concerns theta derivatives and inert-prime divisibility. It does not, in the supplied four-page text, state a bound for a coefficient-matroid minor. Steps~3 and~5 (p.~3) motivate trace and square-root-cancellation arguments; P-02, P-04 and P-06 implement the corresponding curve/sheaf constructions with their exact degree and conductor budgets. Steps~4 and~7 (pp.~3--4) concern determinant and Gan--Gross--Prasad period positivity for specified groups; R-03 identifies why that branching problem cannot be substituted for the tensor subgroup without a proved functorial bridge. Step~8 (p.~4) concerns theta-lift multiplicities; I-03 and I-14 identify the different coefficient-span quantifiers a proposed isolation transfer must preserve.

Step~13 (p.~4) invokes a Bost slope inequality, and step~17 (p.~4) concerns the displayed arithmetic cohomology groups. The supplied note gives no reusable estimates for the ranks, heights or valuations of the ROABP coefficient maps used here. Their possible use therefore has a concrete prerequisite: construct the relevant arithmetic bundle or cohomology object, identify the desired minor or quotient with its invariant, and retain effective size bounds. This is a dependency statement about the supplied inputs, not a claim that such a construction is impossible. The integral and special-fiber analyses below identify where that next construction would act.
